\documentclass{article}

\usepackage{amsmath,amssymb}
\usepackage{xurl}
\usepackage[hidelinks]{hyperref}
\setlength{\emergencystretch}{2em}

% Shared commands for notes in docs/paper-gaps/.

\DeclareUnicodeCharacter{2080}{\ifmmode{}_{0}\else\textsubscript{0}\fi}
\DeclareUnicodeCharacter{2081}{\ifmmode{}_{1}\else\textsubscript{1}\fi}
\DeclareUnicodeCharacter{2082}{\ifmmode{}_{2}\else\textsubscript{2}\fi}
\DeclareUnicodeCharacter{2099}{\ifmmode{}_{n}\else\textsubscript{n}\fi}

\newcommand{\Lean}{\textsc{Lean}}
\ExplSyntaxOn
\NewDocumentCommand{\leanid}{m}{
  \ifmmode
    \textsf{\detokenize{#1}}
  \else
    \group_begin:
    \urlstyle{sf}
    \tl_set:Nn \l_tmpa_tl {#1}
    \tl_replace_all:Nnn \l_tmpa_tl {\_} {_}
    \exp_args:NV \path \l_tmpa_tl
    \group_end:
  \fi
}
\ExplSyntaxOff
\providecommand{\tr}{\operatorname{tr}}
\newcommand{\tnleanIssueBase}{https://github.com/LionSR/TNLean/issues/}
\newcommand{\tnleanPrBase}{https://github.com/LionSR/TNLean/pull/}
\newcommand{\tnleanDocBase}{https://sirui-lu.com/TNLean/docs/}
\newcommand{\ghissue}[1]{\href{\tnleanIssueBase#1}{GitHub issue \##1}}
\newcommand{\ghpr}[1]{\href{\tnleanPrBase#1}{PR \##1}}
\newcommand{\doclink}[1]{\href{\tnleanDocBase#1.html}{\path{#1}}}

% Dirac notation and the identity operator, as in blueprint/src/macros/common.tex.
% \providecommand keeps note-local definitions authoritative.
\usepackage{dsfont}
\providecommand{\ket}[1]{|#1\rangle}
\providecommand{\bra}[1]{\langle #1|}
\providecommand{\braket}[2]{\langle #1 | #2 \rangle}
\providecommand{\ketbra}[2]{|#1\rangle\!\langle #2|}
\providecommand{\Id}{\mathds{1}}
\providecommand{\one}{\mathds{1}}
\providecommand{\C}{\mathbb{C}}
\providecommand{\MN}[1]{M_{#1}(\C)}
\providecommand{\MD}{M_D(\C)}

% External referencing for paper sources.  The build helper writes sanitized
% auxiliary files under build/paper-aux/ and excludes duplicated source labels.
\usepackage{xr}

% Import a generated paper auxiliary file with a caller-supplied label prefix.
% The candidate paths support builds from docs/paper-gaps/, the repository root,
% and build/paper-aux/ itself.
\newcommand{\importpaperaux}[2]{%
  \IfFileExists{../../build/paper-aux/#2.aux}{%
    \externaldocument[#1]{../../build/paper-aux/#2}%
  }{%
    \IfFileExists{build/paper-aux/#2.aux}{%
      \externaldocument[#1]{build/paper-aux/#2}%
    }{%
      \IfFileExists{#2.aux}{%
        \externaldocument[#1]{#2}%
      }{}%
    }%
  }%
}

% General external reference macros
\newcommand{\paperref}[2]{\ref{#1-#2}}
\newcommand{\papereqref}[2]{(\ref{#1-#2})}

% CPSV16 source (1606.00608 / MPDO-22-12-17-2.tex) external document and shortcuts
\importpaperaux{cpsv16-}{1606.00608/MPDO-22-12-17-2-tnlean}
\newcommand{\cpsvref}[1]{\paperref{cpsv16}{#1}}
\newcommand{\cpsveqref}[1]{\papereqref{cpsv16}{#1}}

% Verdict marker: \gapnote{<kind>}{<status>}. Kind is the policy
% classification (clarification, local-correction, scope-restriction,
% unfaithful, false-source, open-gap); status is open, wip (elimination
% underway), resolved, or historical. Machine-read by the site index;
% prints a verdict line.
\newcommand{\gapnote}[2]{\par\noindent\textbf{Verdict:} #1 (#2).\par}


\title{Wolf Equation~(8.103) --- Small-Power Exponent Scope}
\author{}
\date{\today}

\begin{document}
\maketitle

\section{The source statement}

Wolf considers an upper-triangular matrix $X=\Lambda+N\in M_D(\mathbb C)$,
where $\Lambda$ is diagonal and $N$ is strictly upper triangular.  For an
arbitrary submultiplicative norm satisfying $\lVert\Lambda\rVert\le 1$, the
notes state for every $n\in\mathbb N$ that
\begin{align}
  \lVert X^n\rVert
  \le \lVert\Lambda^n\rVert
  +(D-1)n^{D-1}\lVert\Lambda\rVert^{n-D+1}
    \max\{\lVert N\rVert,\lVert N\rVert^{D-1}\}.
\end{align}
The local source is
\path{Notes/WolfNoteTexSource/ch08_distance_measures.tex}, lines~1189--1210,
especially Equation~(8.103) and its derivation from Equation~(8.105).

\section{The exponent defect}

When $n<D-1$, the exponent $n-D+1$ is negative.  The source gives no
convention for this negative power.  In particular, if $\Lambda=0$, then the
factor $\lVert\Lambda\rVert^{n-D+1}$ is not defined in the usual real-power
interpretation.  Replacing the exponent by truncated natural subtraction
would silently change the displayed statement.

The derivation from Equation~(8.105) also factors
\begin{align}
  \lVert\Lambda\rVert^{n-k}
  =\lVert\Lambda\rVert^{n-(D-1)}
   \lVert\Lambda\rVert^{(D-1)-k},
\end{align}
which uses natural exponents only when $D-1\le n$ (and $k\le D-1$).
Thus the proof printed in the notes establishes the displayed natural-power
bound in that range, but does not supply an all-$n$ interpretation.

\section{Formalization scope}

Equation~(8.105), whose exponents satisfy $k\le n$, is formalized for every
$n\in\mathbb N$.  Equation~(8.103) is formalized with the explicit hypothesis
$D-1\le n$.  Its refined constant $(D-1)\binom{n}{D-1}$ is formalized under
the source's stated hypothesis $2(D-1)\le n$, which implies the first range
condition.  Both results retain the arbitrary submultiplicative seminorm
parameter; no particular matrix norm replaces it.  The corresponding Lean
declarations are:
\begin{itemize}
  \item \leanid{Matrix.wolf_eq_105_seminorm} for the unrestricted
    Equation~(8.105) word bound;
  \item \leanid{Matrix.wolf_eq_103} for the natural-power form under
    $D-1\le n$;
  \item \leanid{Matrix.wolf_eq_103_refined} for the refined constant under
    $2(D-1)\le n$.
\end{itemize}

\section{Verdict}

\textbf{Verdict: source scope restriction.}

This is not a missing Lean proof.  The
formal statement records exactly the range in which the printed
natural-exponent derivation is meaningful.  Eliminating the restriction would
require Wolf to specify an interpretation for negative powers, including the
case $\lVert\Lambda\rVert=0$.  Tracking: \ghissue{5837}.

\bibliographystyle{alpha}
\bibliography{references}

\end{document}
